% ============================================================
% Secao 1: Introducao
% ============================================================

\section{Introdu\c{c}\~ao}
\label{sec:introducao}

Modelos de linguagem de grande escala (LLMs) alcan\c{c}aram capacidades
not\'aveis em gera\c{c}\~ao de texto, racioc\'inio e seguimento de instru\c{c}\~oes
\cite{brown2020gpt3, touvron2023llama}. Contudo, estes modelos sofrem de uma
limita\c{c}\~ao fundamental: \textbf{n\~ao sabem quando n\~ao sabem}.

A sa\'ida de um LLM convencional \'e um vetor de logits sobre o vocabul\'ario,
seguido de uma opera\c{c}\~ao softmax que produz probabilidades. Embora a entropia
desta distribui\c{c}\~ao possa servir como proxy grosseiro de incerteza, ela
confunde incerteza aleat\'oria (inerente aos dados) com incerteza epist\^emica
(devida \`a falta de conhecimento do modelo). Al\'em disso, modelos treinados com
\textit{cross-entropy} tendem a ser \textbf{mal calibrados}: a confian\c{c}a
reportada n\~ao corresponde \`a acur\'acia real \cite{guo2017calibration}.

Trabalhos anteriores abordam este problema de forma extr\'inseca:
ensembles \cite{lakshminarayanan2017simple}, Monte Carlo Dropout
\cite{gal2016dropout}, ou sistemas de orquestra\c{c}\~ao que analisam respostas
ap\'os a gera\c{c}\~ao. Estas abordagens adicionam lat\^encia, custos computacionais
e n\~ao permitem que o sinal epist\^emico influencie o pr\'oprio processo de
gera\c{c}\~ao.

\subsection{Contribui\c{c}\~oes}

Neste trabalho, propomos o \aletheion{}, um LLM \textit{decoder-only} que integra
o sistema epist\^emico \textbf{intrinsecamente} na arquitetura neural:

\begin{enumerate}
    \item \textbf{Tomografia Epist\^emica por Token}: Cada token gera, al\'em de
    logits, um conjunto de 30+ campos epist\^emicos organizados em 3 tiers,
    incluindo incerteza decomposada ($q_1$, $q_2$), confian\c{c}a calibrada,
    coordenadas em manifold 5D e estado cognitivo.

    \item \textbf{Directional Relational Manifold (DRM)}: Cada token \'e
    mapeado para um espa\c{c}o 5D com geometria Riemanniana aprendida, onde
    dist\^ancias geod\'esicas definem a confian\c{c}a epist\^emica via
    \textit{Metric-Aware Distance} (MAD).

    \item \textbf{Vetor de Intencionalidade (VI)}: Um sistema de monitoramento
    e corre\c{c}\~ao da sa\'ude do manifold, baseado em 4 componentes ($\phi$)
    que detectam colapso dimensional e aplicam corre\c{c}\~oes determinisicas.

    \item \textbf{Navegador MPC}: Controle preditivo no manifold com
    \textit{beam search} (K=4, D=3) que planeja trajet\'orias \'otimas sobre
    12 a\c{c}\~oes de interven\c{c}\~ao.

    \item \textbf{9 M\'odulos de Extens\~ao}: EidosDecay, Filosofia3 (conflito
    $\phi$-$\psi$), SelfModel, Grounding, PlasticityGate, MPL/Frontier,
    MOPsi, CausalState e Metacognitive --- totalizando ~364K par\^ametros
    (<0.3\% do modelo).

    \item \textbf{Loss Composta com Annealing}: 14 componentes de perda com
    \textit{warmup} progressivo que estabiliza o treinamento.

    \item \textbf{Continual Learning}: EWC com acumula\c{c}\~ao online de Fisher
    e \textit{Experience Replay} com reservoir sampling integrados no
    \textit{trainer}.
\end{enumerate}

\subsection{Origem e Autoria Intelectual}

Os conceitos epist\^emicos implementados no \aletheion{} --- incluindo o
\textbf{Differential Reference Manifold (DRM)}, \textbf{Metric-Aware Distance (MAD)},
\textbf{Vetor de Intencionalidade (VI)}, \textbf{Navegador MPC},
\textbf{MOPsi} (m\'odulo orientador de $\psi$), \textbf{MPL} (projetor de longo prazo),
\textbf{EidosDecay}, \textbf{Filosofia3} (conflito $\phi$-$\psi$),
\textbf{SelfModel} (consci\^encia), \textbf{Termodin\^amica Computacional}
(custo de Landauer, plasticidade, din\^amicas irrevers\'iveis) e o
\textit{framework} de \textbf{Tomografia Epist\^emica} --- foram
\textbf{originalmente concebidos, projetados e implementados por Felipe Maya Muniz}
no projeto ATIC (Adaptive Turing Intelligent Cognition, 2025--2026), onde operam
como m\'odulos Python num pipeline de orquestra\c{c}\~ao.

No ATIC, estes sistemas calculam m\'etricas epist\^emicas \textit{post-hoc}
sobre respostas de LLMs. A contribui\c{c}\~ao central do \aletheion{} \'e a
\textbf{transposi\c{c}\~ao destes conceitos para componentes \texttt{nn.Module}
trein\'aveis} dentro da pr\'opria rede neural, permitindo aprendizado
\textit{end-to-end} de sinais epist\^emicos que influenciam diretamente a
gera\c{c}\~ao de tokens.

\subsection{Motiva\c{c}\~ao Filos\'ofica}

O nome \textit{Aletheion} deriva do grego $\alpha\lambda\acute{\eta}\theta\epsilon\iota\alpha$
(\textit{al\=etheia}), que significa ``desvelamento'' ou ``verdade'' no sentido
heideggeriano. A proposta \'e que o modelo n\~ao apenas gere texto, mas
\textbf{desvele} o seu pr\'oprio estado epist\^emico --- tornando expl\'icito
o que sabe, o que n\~ao sabe, e o qu\~ao confi\'avel \'e cada token.

\subsection{Organiza\c{c}\~ao do Paper}

A \Cref{sec:arquitetura} descreve a arquitetura geral do transformer.
A \Cref{sec:epistemico} detalha o \texttt{EpistemicHead} e seus sub-heads.
As se\c{c}\~oes \ref{sec:drm}--\ref{sec:mpc} aprofundam os componentes centrais
(DRM, MAD, VI, MPC). A \Cref{sec:extensoes} cobre os 9 m\'odulos de extens\~ao.
A \Cref{sec:loss} apresenta a loss composta. A \Cref{sec:continual} discute
continual learning. A \Cref{sec:escalabilidade} aborda escalabilidade.
A \Cref{sec:conclusao} conclui.
